\section{Expressions over the original alphabet}
\label{sec:compilation}

Both split interfaces have been proved, including their universal
soundness under independent guesses. We have not yet assigned a height
to any of their component languages. Finite-state decidability of a test
would establish regularity, but would not give the height-one expressions
that the split identity needs.

The relevant restriction comes from each component's own end rule. A
marker return gives a bounded end. Otherwise a fixed periodic seed is
continued until its first mismatch, followed by a bounded tail. Hence
there is only one unbounded repeated word, and finitely many choices of
its fixed surrounding words. We first give the exact compilation lemma
for that situation, including additional end comparisons made at false
origins. This develops the template method of
\cite[Lemma~8.1 and Proposition~8.2]{HeightThirteen} and
\cite[Lemma~7.1 and Proposition~7.2]{HeightFour} for the present guards.

\subsection{Periodic templates and exact-end checks}

\begin{lemma}[Compilation on periodic templates]\label{compile:templates}
Fix a finite alphabet $\Sigma$, a finite monoid $M$, and a morphism
$T:\Sigma^*\to M$. Consider finitely many templates
\begin{equation}\label{compile:eq-template}
                     a c^i b,\qquad i\ge0,
\end{equation}
where $a,b,c\in\Sigma^*$ are fixed for each template and $c\ne\eps$.
On each template suppose acceptance is determined by a fixed finite
calculation using the following data:
\begin{enumerate}
\item letters in finitely many bounded windows at fixed integer offsets
from the argument's beginning or end;
\item order and availability of positions selected from fixed finite sets
of such offsets, with unavailable required data causing rejection;
\item monoid products on available ordered intervals between these positions,
and lengths or differences modulo fixed positive integers;
\item agreement of an interval between these positions with one of finitely
many fixed periodic continuations, together with any required mismatch
letter at its endpoint.
\end{enumerate}
Finite choices, universal tests, and uniqueness tests on fixed finite lists
are allowed in the calculation. Then the accepted words have a generalized
expression of height at most one over $\Sigma$. The same conclusion holds
after adding finitely many exceptional words.
\end{lemma}

\begin{proof}
Fix a template $ac^ib$. For sufficiently large $i$, the relative order
and availability of every pair of bounded-offset positions is fixed.
Adding another full copy of $c$ then leaves a bounded window near either
end unchanged. Products between endpoints near the same end are likewise
constant. Every interval product crossing the growing middle has form
\[
                       u\,T(c)^{i-i_0}\,v
\]
for fixed $u,v\in M$ and an integer $i_0$, after taking sufficiently
large $i$. Powers in a finite monoid are ultimately periodic. Thus all
these products, and all fixed-modulus length data, are ultimately
periodic functions of $i$.

For a comparison with a periodic continuation of period $d_0$, separate
the finitely many classes of $i$ modulo $d_0$. On a class, the phase of
every bounded endpoint is fixed. Comparisons in bounded pieces therefore
stabilize. A growing middle compares the repetition of $c$ with the
specified continuation in a fixed phase. If they disagree, a disagreement
occurs within $|c|d_0$ consecutive positions: this integer is a common
period of the two two-sided repetitions. Every sufficiently long middle
then fails the comparison. If there is no disagreement in such a block,
the entire growing comparison agrees, and only the fixed beginning and
end comparisons remain. A prescribed first mismatch is exactly an
agreement test up to a specified endpoint and a differing letter there,
so it is covered too. Bounded comparison intervals reduce to the earlier
window case.

There are only finitely many data and phase cases. Choose a threshold
and a common period beyond which their whole array is periodic. Any
fixed calculation on that finite array, including a universal check or
an exactly-one condition, has an ultimately periodic set of accepted
exponents. Such a set is a finite set together with finitely many
progressions $i_1+\ell\mathbb N$, $\ell>0$. A progression is expressed by
\[
                       ac^{i_1}(c^\ell)^*b.
\]
All its fixed powers are words over $\Sigma$, so this has height one.
A finite exceptional set contributes finitely many words. Taking finite
unions over the progressions and templates proves the assertion on the
original alphabet.
\end{proof}

An exact end check at a fixed candidate origin provides precisely the
template cover required by this lemma. If the smaller end search $\eta$
returns a marker, that marker has a bounded offset and the end is bounded.
Otherwise its seed is a fixed-length word at a fixed position, with a
short-period continuation. Let $a$ be the bounded word from the argument's
own beginning to the seed's right endpoint. Make finite choices of $a$,
the continuation's period and phase, and the tail of length
$L_{\rm tail}$ starting at the first mismatch. If the mismatch follows
$i$ full copies of the period word and one partial copy, put the fixed
partial copy and tail into $b$ in \eqref{compile:eq-template}. Require the
first letter of the tail to differ from the periodic continuation. These
choices give a finite template cover, with bounded cases handled separately.

This is a statement about the suffix language itself. Its proposed origin
is one of finitely many offsets from its own beginning, and its own end
check yields the cover whether or not that origin could belong to a
successful paired split. There is no ambient-episode assumption here.

Multiple end guards also cause no additional unbounded search. Every
guard that uses a mismatch requires that mismatch at the proposed
argument end minus $L_{\rm tail}$. After choosing its seed and period in
finite cases, it tests agreement up to that position and disagreement
there. The reference guard and decoded-origin guard are therefore both
among the comparisons of Lemma~\ref{compile:templates}. A missing seed
or window rejects the test, rather than making the template argument
conditional on inaccessible letters.

\subsection{The individual component languages}

\begin{proposition}\label{compile:components}
The episode language $E$, its skipped sublanguage $D$, and every suffix
language in the outer and inner splits have generalized star height at
most one over $\Sigma$. Every prefix language in these splits has height
zero.
\end{proposition}

\begin{proof}
All prefix lengths belong to fixed finite schedules. Each prefix language
is therefore finite and has height zero.

For the other components, first obtain a finite template cover from an
exact end check. The languages $E$ and $D$ use their episode end rule.
An outer suffix uses the reference-origin end check of its common guard.
An inner suffix always checks the decoded origin's end and, when a larger
search is absent, also checks its reference origin's end. Separate the
finitely many possible decoded origins into cases. The preceding template
construction applies in every case, including standalone suffixes with
incorrect hypotheses or inconsistent tag fields.

It remains to verify the allowed data on these covers. Every schedule,
origin hypothesis, translated origin in a stability check, and marker
inspection radius is fixed and finite. The integers in a bounded shift
range form a finite list, even if that list is very large. A bounded marker
search therefore selects a position from finitely many offsets from the
argument's beginning, or returns absence. All comparisons of such searches,
their first-marker choices, and all window-availability requirements are
finite calculations from the bounded data.

The clock anchor $Q_\tau$ is either one of those bounded larger-search
returns or the mismatch anchor supplied by an exact end check. In the
latter case it is the argument end minus $L_{\rm tail}$. The residue
$(Q_\tau-\tau)\bmod B$ and the clock value $\varphi(Q_\tau-\tau)$ use only
fixed moduli. Consequently the predicates $Z$ and $A_*$, and the universal
and uniqueness checks in the decoder, use the data of
Lemma~\ref{compile:templates}. Where an absent larger search requires a
reference end check, its periodic agreement requirement is covered by
the last part of that lemma. The additional decoded end check is covered
in the same way; there is no need to search past the proposed endpoint.

The actual product intervals in the outer tests have endpoints at selected
bounded search positions, cuts, or the prescribed end. The inner stencil
boundaries have fixed offsets under each origin hypothesis. Computing
$\psi_e(m')$ for every $m'\in M$ uses products from its final stencil
boundary to selected later boundaries or to the end, together with the
same bounded marker and clock data. The defining formula accepts every
hypothetical $m'$; it adds no test that a prefix realizing $m'$ exists.
All shift-table certificates range over fixed finite monoids and vector
spaces. Hence they are finite calculations from the listed data. In
particular, no component calls a graph-language test for a sequence of
episodes.

We have checked every allowed source of data, not just regularity of the
resulting test. Lemma~\ref{compile:templates} therefore gives the claimed
height-one expressions for these standalone languages. The finite tuple
of tags, states, and residues merely indexes finite operations on those
expressions. Every word appearing inside a template, including $c^\ell$
under its single star, is a word over $\Sigma$. No height-preserving
alphabet substitution is being assumed.
\end{proof}

\subsection{Words with no complete episode prefix}

Define the final remainder language
\[
              N_{\rm end}=\neg(E\,\neg0).
\]
Since $\neg0=\Sigma^*$ has height zero, this language has height at most
one by Proposition~\ref{compile:components}. We also need to retain the
monoid product on a remainder.

\begin{lemma}\label{compile:remainder}
For every $m_0\in M$, the language
$N_{\rm end}\cap T^{-1}(m_0)$ has height at most one over $\Sigma$.
Every word has a unique factorization into complete episodes followed by
a word in $N_{\rm end}$.
\end{lemma}

\begin{proof}
Choose a fixed length threshold beyond every bounded decision window for
an episode beginning at zero and beyond every possible marker-present
episode end. A longer word in $N_{\rm end}$ cannot have a present smaller
end search: its prescribed complete episode would already be a prefix.
Thus its available seed has a short-period continuation. Either this
continuation has no mismatch before the word ends, or the first mismatch
$q_0$ satisfies
\[
                       |w|-q_0<L_{\rm tail}.
\]
Otherwise the prefix ending at $q_0+L_{\rm tail}$ would be a complete
episode. The margin construction ensures that all its required bounded
data are internal even when the mismatch is as early as the rule allows.

In the first case, finite choices of the beginning through the seed,
period, phase, and final partial period give templates
\eqref{compile:eq-template}. In the second, append a portion of length
less than $L_{\rm tail}$ starting at the mismatch, chosen from a finite
set. These too are such templates. Shorter words form a finite language.
On each template, the monoid-value test is ultimately periodic by
Lemma~\ref{compile:templates}. Construct the height-one expression for
that value on the template cover and intersect it with $N_{\rm end}$.
The result is exactly $N_{\rm end}\cap T^{-1}(m_0)$, with the same bound.

Repeatedly remove an episode prefix whenever one exists. Episodes are
nonempty, so the procedure terminates, leaving a word in $N_{\rm end}$.
The prefix-code property makes the removed prefix unique at every step.
No alternative factorization can stop sooner, since its remainder would
still have an episode prefix. This proves the factorization assertion.
\end{proof}

\subsection{Two splits and the main theorem}

\begin{proof}[Proof of Theorem~\ref{thm:main}]
The empty alphabet was handled at the start of the construction. For a
nonempty alphabet, Proposition~\ref{compile:components} supplies all
component bounds required by Lemma~\ref{alg:split}.

First use the inner split of Proposition~\ref{inner:split} on $D$, with
no skipped episodes. The graph languages of its empty skip domain are
the identity tests on $\{\eps\}$ and have height zero. The split identity
therefore gives height at most two for every graph language of $G_e$ on
$D^*$. Its linear part is $\widehat F_e$. Boolean affine recovery gives
graph tests for products of these linear parts at the same height, and
evaluation on $(v,1)$ recovers every graph language of $F_e$ on $D^*$.

Use these as the skip graph languages for the outer split of
Proposition~\ref{outer:split}. They refer to the same deterministic
updates $F_e$ as that split, including on the skipped episodes. A second
application of the split identity yields height at most three on $E^*$:
\[
                  0\ \longrightarrow\ 2\ \longrightarrow\ 3.
\]
Boolean recovery now removes the affine shifts of $F_e$. Their linear
parts are the right actions of the original episode products. The image
of the identity basis vector distinguishes each element of $M$, so we
obtain height-at-most-three tests
\[
                       A_d=E^*\cap T^{-1}(d),\qquad d\in M.
\]
Let $B_f=N_{\rm end}\cap T^{-1}(f)$ be the height-one tests of
Lemma~\ref{compile:remainder}. If $T$ recognizes the regular language
with accepting subset $J\subseteq M$, the expression
\[
                \bigcup_{\substack{d,f\in M\\df\in J}} A_d B_f
\]
defines it. Existence of the episode--remainder factorization gives
completeness, and every accepted factorization has the stated product
$df$, giving soundness. Concatenation and finite union preserve the
maximum height, which is three. Every expression used here is over the
given alphabet $\Sigma$, with complement relative to $\Sigma^*$.
\end{proof}
